% ============================================================
% Model — Paper A (JNCA): the organizational communication graph in time
% and the cascade invariant. Variables only; numbers live in the Case Study.
% ============================================================
\section{Model}
\label{sec:model}

The model formalizes the organization's communication graph in time and defines the invariant that a
compromised account violates. All quantities here are variables; the numbers that instantiate them are stated
in the Case Study (Sec.~\ref{sec:results}). Figure~\ref{fig:model} gives an overview.

\begin{figure}[pos=htbp]\centering
\resizebox{0.98\textwidth}{!}{%
\begin{tikzpicture}[>=Latex, font=\small, node distance=6mm,
  box/.style={draw, rounded corners=2pt, align=center, minimum height=8mm, inner sep=3pt, fill=blue!4},
  lay/.style={draw, align=center, minimum width=22mm, minimum height=6mm, inner sep=2pt, fill=blue!8},
  det/.style={box, fill=red!5}, flow/.style={-{Latex[length=2mm]}, thick}]
  \node[box] (org) {Organization\\\scriptsize accounts $V$, roles, partners};
  \node[lay, right=14mm of org, yshift=5mm] (l1) {contacts $E$\\\scriptsize who writes to whom};
  \node[lay, below=2mm of l1] (l2) {rhythm $\tau$\\\scriptsize active windows};
  \node[box, right=16mm of l1, yshift=-5mm] (inv) {Cascade invariant\\\scriptsize off-rhythm $tc{=}0$ $+$\\receiving-side wave $m_v^{(b)}$};
  \node[det, right=14mm of inv] (d) {temporal detector\\\scriptsize node memory};
  \node[box, right=12mm of d, fill=green!6] (eff) {Effectiveness\\\scriptsize leak-aware, low-FPR,\\base rate (Sec.~\ref{sec:effectiveness})};
  \node[draw, dashed, rounded corners, fit=(l1)(l2), inner sep=4pt, label=above:{\footnotesize communication graph $G$ (Def.~\protect\ref{def:cg})}] (cg) {};
  \draw[flow] (org) -- (cg);
  \draw[flow] (cg.east) -- (inv.west);
  \draw[flow] (inv) -- (d);
  \draw[flow] (d) -- (eff);
\end{tikzpicture}}
\caption[General model]{The model, general to detailed (left to right). An organization induces the
communication graph $G$ with a rhythm per account; a compromised account violates it by writing off-rhythm
and by propagating, which leaves a receiving-side wave; a temporal detector with node memory reads that wave,
and is judged by one \emph{definition of effectiveness}.}
\label{fig:model}
\end{figure}

\begin{table}[pos=tbp]\centering\small
\caption{Notation used throughout the paper. Concrete instantiations of the free
parameters ($N$, $T$, $p_\mathrm{inf}$, $H$) are given in the Case Study (Sec.~\ref{sec:results}).}
\label{tab:notation}
\begin{tabular}{ll}
\toprule
\textbf{Symbol} & \textbf{Meaning} \\
\midrule
$G=(V,E,\tau)$ & organizational communication graph with rhythm \\
$V$ & set of accounts (people/roles), $|V|=N$ \\
$E$ & contacts: $(s,v)\in E$ iff $s$ corresponds with $v$ \\
$\tau(s)$ & temporal rhythm (active windows) of sender $s$ \\
$T$ & number of time buckets partitioning the period \\
$b$ & a single time bucket, $b\in\{1,\dots,T\}$ \\
$tc(s,v,t)$ & temporal consistency of $(s,v)$ at time $t$ \\
$s,v$ & sender, victim (recipient); $s_0$ the compromised account \\
$p_\mathrm{inf}$ & per-contact infection probability of a cascade \\
$H$ & maximum cascade depth (hops) \\
$K$ & number of carriers of a stealthy attacker (one send per bucket) \\
$Z$ & static-GNN node embedding matrix, Eq.~\eqref{eq:static} \\
$d_v^{(b)}$ & in-degree of node $v$ within bucket $b$, Eq.~\eqref{eq:temp-mem} \\
$m_v^{(b)}$ & temporal-GNN memory of node $v$ after bucket $b$ \\
\bottomrule
\end{tabular}
\end{table}

\subsection{The organizational communication graph}
\label{sec:cg}
An organization's communication is relational and rhythmic: who corresponds with whom follows the
organizational chart (supervisors, coworkers, business partners), and each account is active in
characteristic windows of the week. Both facts are observable from mail headers alone, without content.

\begin{definition}[Organizational communication graph]
\label{def:cg}
A communication graph is $G=(V,E,\tau)$, where $V$ is the set of accounts (people/roles), $E\subseteq V\times V$
is the set of contacts ($(s,v)\in E$ iff $s$ corresponds with $v$; induced by the organization's hierarchy,
coworker, and partner relations), and $\tau:V\to\mathcal{T}$ assigns to every sender a \emph{temporal
rhythm}, the set of its active time windows. The temporal consistency of a message from $s$ to $v$ at time
$t$ is $tc(s,v,t)=\mathbf{1}[t\in\tau(s)]$.
\end{definition}

\noindent A message from a known contact in its rhythm, $(s,v)\in E$ and $tc=1$, is the normal case. A known
contact writing \emph{off-rhythm}, $tc=0$, is the canonical signature of a compromise: the relation is
genuine, the timing is not. This single bit is weak on its own (legitimate mail also arrives off-hours, and an
attacker with mailbox access can mimic the rhythm), which is why the invariant below reads the dynamics of the
whole neighborhood rather than one message.

\subsection{Threat model: the compromised account}
\label{sec:threat}
The attacker has gained mailbox access to an internal account $s_0$ (the \emph{patient zero}) and wants to
deliver to victims messages classified as legitimate. From $s_0$ the attacker mails the account's own
contacts; each recipient becomes a carrier with probability $p_\mathrm{inf}$ and propagates to its own
contacts, up to depth $H$. The attack is therefore a \emph{cascade} on $G$ (Fig.~\ref{fig:cascade} in
Sec.~\ref{sec:method}). The term follows \citet{ho2019detecting}, who coined \emph{lateral phishing} for
phishing sent from compromised enterprise accounts and characterized it at scale; the multi-hop cascade with
a per-contact infection probability is our formalization of the compounding they observed, in which newly
compromised recipients become the next senders. Every single message in it is locally plausible: the sender is a genuine contact,
the address is real, and the content can be clean.

On capabilities: having mailbox access, the attacker knows the victim's rhythm and can mimic it; the attacker
can also \emph{pace} the sends---$K$ carriers, one send per time bucket---so that no sending burst ever appears
(the \emph{stealthy} attacker, who defeats volumetric detection by construction); and the attacker can restrict
the attack to a single recipient with no propagation at all (the \emph{surgical} attack). Out of scope is an
adaptive adversary with full knowledge of the graph and training poisoning, and external impersonation from
outside the graph, which requires knowledge beyond $G$ and is treated in a separate study. We treat the signal
as a cost increase for the attacker; no detection is final.

\subsection{The cascade invariant}
\label{sec:invariant}
Lateral phishing is a propagation process, and propagation has two sides. On the \emph{sending} side a
compromise shows up as a surge in one account's volume, or as an unusual path to a rare target; this is what
practical defenses read, and it is what a patient attacker controls. On the \emph{receiving} side the cascade
leaves a signature that no single message reveals and that the attacker cannot pace away: within the wave's
window many of $v$'s neighbors receive mail, and the sender $u$ of the message to $v$ was itself hit just
before. This receiving-side wave is the invariant. A learned temporal model with node memory registers it by
scoring an event using the state from \emph{before} its bucket, so that it can only use the past (made precise
in Sec.~\ref{sec:method}); a static model of the same topology, which sees who is connected to whom but not
in what order, cannot. The invariant is testable by ablation: shuffling the time order of events destroys
the order it relies on (together with within-bucket co-occurrence, Sec.~\ref{sec:m-detectors}), so the drop in
performance under a time shuffle bounds how much of the signal is dynamics rather than topology.

\subsection{Definition of detection effectiveness}
\label{sec:effectiveness}
Every result in the Case Study follows a standard that is detector-independent and operational; the
Methodology (Sec.~\ref{sec:method}) is the protocol that operationalizes it.

A high benchmark accuracy alone does not attest to detecting lateral phishing. A method works operationally
for the targeted tail only when its advantage is temporal, vanishing under a control that destroys the
temporal order, and transferable, so that the qualitative hierarchy of methods---which class wins, and where it
breaks down---carries over to other topologies. Effectiveness is judged by sensitivity at a low
false-positive rate with explicit base-rate accounting, rather than by headline AUC. The leak-aware protocol of
Sec.~\ref{sec:method} implements this standard.
